\section{Core charges and a low-dimension cutoff}
\label{rec:sec-cutoff}

The previous recursions used either the first-row level or a changing
tensor alphabet.  Here the cost of a restriction is charged directly to
boxes lying beyond the first several rows and columns of a Young diagram.
The goal is again a negative dimension power for one sweep, but the
intermediate statement concerns a specific projection at its midpoint:
types with small product dimension have exponentially small overlap with
the full sweep.

There are two bounds for each weighted cycle factor.  One depends only on
the diagrams, and the other depends on their actual positive traces.
We choose between them while expanding the recursion.  The second choice
must use the traces at the current frontier; a later upper bound cannot
replace them before the choice is made.  We make that invariant explicit
below.

Logarithms in this section are to base two, unless $\ln$ is written.
For a nontrivial diagram $\lambda\vdash n$, put
$F_\lambda=\log_2D_\lambda$, $k=n-\lambda_1$, and
$L=\log_2(n/k)$.  For a box $x$ let
$a(x)=\min(\operatorname{row}(x),\operatorname{column}(x))$, with
indices beginning at one.  We call a type surviving if its sweep matrix
is nonzero; for $n\ge2$ it then has at most $n/2$ rows.  Traces with zero contribution
may be omitted throughout.  For $s\ge0$ define the core count
\[
 u_s(\lambda)=\#\{x\in\lambda:a(x)>s\}.
\]
Thus the core is what remains after deleting the first $s$ rows and the
first $s$ columns.  Although its box coordinates shift under deletion,
its remaining row lengths form a straight Young diagram.

\begin{theorem}\label{rec:core-power-main}
There is an absolute $c>0$ such that
$\|T_n|_{V_\lambda}\|_{\op}\le2^{-cF_\lambda}$ for every
surviving nontrivial type.
\end{theorem}

\subsection{Dimension charges and the local sparse estimate}
We need two elementary inputs, uniformly over surviving nontrivial types
for all sufficiently large dyadic $n$:
\begin{equation}\label{rec:core-dimension-sparse}
 F_\lambda\ge b k,\qquad
 1\le k\le n^{1/100}\ \Longrightarrow\
 \|T_n|_{V_\lambda}\|\le2^{-b_1k\log_2n},
\end{equation}
where $b,b_1>0$ are absolute.  We prove the sparse estimate locally
because its range grows with $n$.

For the dimension estimate, write $\tau$ for the diagram below the first
row.  The hook formula gives
\[
 D_\lambda=\binom nkD_\tau
 \bigg/\prod_{j\le\lambda_1}
       \left(1+\frac{\tau'_j}{\lambda_1-j+1}\right).
\]
The denominator is at most $2^k$: for an integer $t\ge0$ and
positive integer $a$, $1+t/a\le2^t$, and
$\sum_j\tau'_j=k$.  When $k\le n/8$, the bound
$\binom nk\ge(n/k)^k$ proves $F_\lambda\ge2k$.
For $k\ge n/8$, first suppose that a first row or first column has
length $a\ge n/10$.  Transpose if needed and retain that row together
with $v=\lfloor n/100\rfloor$ boxes below it.  Such a subdiagram
exists: after transposition its tail has at least $n/8$ boxes if the
chosen row was the original first row, and at least $n/2$ if it was
the first column.  Tableaux of a subdiagram extend to the full diagram.
The same formula on this smaller diagram gives
$D_\lambda\ge\binom{a+v}{v}2^{-v}$, exponential in $n$.
If both the first row and first column are shorter than $n/10$, every
hook is shorter than $n/5$ and Stirling's inequality gives
$D_\lambda\ge(5/e)^n$.  These cases prove the first estimate.
They also show $F_\lambda=O(k\log n)$ on the sparse range, the
upper bound following from $D_\lambda\le\binom nkD_\tau$ and
$D_\tau\le\sqrt{k!}$.

To prove the second estimate, realize $V_\lambda$ on ordered injections
of $k$ labels.  A function of that exact position type is harmonic in
each label slot, because the type does not occur on $k-1$ slots.
For distinct input positions $x_a,x_b$, let $h_{ab}$ be their largest
differing bit index in the sweep order.  For $J\subseteq[k]$ put
\[
 g_J(x,y)=\prod_{a<b:\,a,b\in J}
 \left(1-\mathbf1_{\{y_{a,<h_{ab}}=y_{b,<h_{ab}}\}}
                (-1)^{y_{a,h_{ab}}+y_{b,h_{ab}}}\right).
\]
The individual route from an input to an output is unique.  Two labels
can first share a switch only at $h_{ab}$, when their earlier output
bits agree.  Equal output bits also at $h_{ab}$ make the prescription
impossible, giving a zero factor; opposite bits give one common fair
coin in place of two, giving a factor two.  Otherwise the factor is one.
Thus $n^{-|J|}g_J$ is exactly the prescribed subtuple transition
probability, including zero for incompatible routes.

The number of shared switches among $s$ compatible paths is at most
$s\log_2s/2$, by splitting at the last coordinate and adding
$\min(s_0,s_1)$ to the two child counts.  The binary entropy inequality
pays this addition.  Hence $0\le g_J\le k^{k/2}$.
On the exact position type we may replace $g_{[k]}$ by
\[
 g'(x,y)=\sum_{J\subseteq[k]}(-1)^{k-|J|}g_J(x,y):
\]
every omitted-slot term is annihilated by harmonicity.  This sum has
absolute value at most $2^kk^{k/2}$.  Moreover it vanishes if some
label is isolated in the potential-contact graph, whose edge $ab$ is
the event $y_{a,<h_{ab}}=y_{b,<h_{ab}}$; the terms with and without
that isolated label then cancel in pairs.

A graph without isolated vertices contains a spanning forest without
isolated vertices, and such a forest has at least $\lceil k/2\rceil$
edges.  For independent uniform, unrestricted binary strings $x,y$, declare
that $ab$ is not an edge when $x_a=x_b$; on distinct input pairs use
the potential-contact condition above.  An edge constraint in a forest costs $d/n$ when a leaf is integrated out:
the possible last differing input bit $h$ has probability $2^{h-1}/n$,
and agreement of the earlier output bits costs $2^{-(h-1)}$.
Successive leaf removal therefore multiplies these costs.  There are
at most $k^{2k+1}$ forests to include in the union bound.  Conditioning
both endpoint tuples to be injective costs at most two on
$k\le n^{1/100}$ for large $n$.  Consequently
\[
 \Pr(\text{no isolated label})
 \le2k^{2k+1}(d/n)^{\lceil k/2\rceil}
 \le2^{-(2/5)k\log_2n}\qquad(k\ge2).
\]
The tuple kernel density relative to uniform has the additional factor
$(n)_k/n^k\le1$.  Its compressed Hilbert--Schmidt norm is therefore
at most $2^kk^{k/2}$ times the square root of this probability.
This is $2^{-b_1k\log_2n}$ for an absolute $b_1>0$, after increasing
the lower size threshold.  The operator norm is smaller, and at level
one the sweep is exactly zero.  This proves~\eqref{rec:core-dimension-sparse}.

The charge to the deep core has the following form:
\begin{equation}\label{rec:core-charge}
 \sum_{x:\,a(x)\ge R}\{L+\log_2 a(x)\}
 \le(1+o(1))F_\lambda\qquad(R\longrightarrow\infty),
\end{equation}
uniformly over surviving nontrivial diagrams of sufficiently large size.
Here the $o(1)$ depends only on $R$, not on the diagram.
For any straight diagram of size $u$, a box at minimum index $a$ has
hook length at most $2u/a$: each of its row and column lengths is at
most $u/a$.  Apply this to $\tau$ in the first-row hook formula
above.  The coordinates below the first row shift by one, changing
$\log a$ by at most a constant.  Stirling's bound then gives
\[
 \sum_{x\text{ below row }1}\{L+\log_2a(x)\}
 \le F_\lambda+O(k).
\]
If $L>\sqrt{\log_2R}$, the lower bound
$F_\lambda\ge kL-O(k)$ absorbs that error uniformly as $R\to\infty$.

In the complementary case, remove the first
$q=\lfloor\sqrt R\rfloor$ rows and columns.  Let the remaining core
have size $u$.  Applying the additive-error bound just proved to boxes
with $a>q$ gives
$u=O((F_\lambda+k)/\log R)=O(F_\lambda/\log R)$.
The dimension of this core is at most $D_\lambda$, by extending its
tableaux after translating it to the upper left.  Its hook formula
therefore pays $\sum\log_2(a-q)$ up to $O(u)$.
For charged boxes $a\ge R$, replacing $\log_2(a-q)$ by
$\log_2a$ costs $o(1)$ per box, and adding $L$ costs at most
$u\sqrt{\log_2R}=o(F_\lambda)$.  This proves
\eqref{rec:core-charge}.

\subsection{Two bounds for a weighted cycle factor}

The positive network recursion is $K_n=H(K_m\otimes K_m)H$,
where $H$ is the orthogonal average of the outer matching switches.
For a positive $D$ and $p\ge2$,
\begin{equation}\label{rec:cutoff-compressed-jensen}
 \Tr(HDH)^p\le\Tr(HD^{p/2}H)^2.
\end{equation}
Indeed take an eigenbasis of $HDH$ on the range of $H$.
Scalar Jensen gives $\langle v,Dv\rangle^{p/2}
\le\langle v,D^{p/2}v\rangle$ for each basis vector.  Squaring and
summing is bounded by the Hilbert--Schmidt square of the compression.
This proof uses scalar convexity and does not require operator
convexity of the power function.

Resolve the two halves into types $\mu,\theta$.  For the moment,
let $M_\mu$ and $M_\theta$ be arbitrary positive matrices on these
two types, and put $t_\mu=\Tr M_\mu^2$ and
$t_\theta=\Tr M_\theta^2$.  In the unrestricted recursion they
will be $K_m^{p/2}$ on the respective types; allowing arbitrary positive
matrices will also cover the midpoint projections below.
Write $c(g)$ for the number of cycles of a permutation $g$, including
fixed points.  Expansion of the matching projection gives the regular trace bound
\begin{equation}\label{rec:cutoff-weighted-cycle}
 \frac{\Tr(EHEH)}{(2m)!}
 \le4^{-m}\sum_z|e(z)|^2\,2^{c(\alpha^{-1}\zeta)}.
\end{equation}
Here $E$ is the element of the half-preserving group algebra whose only
nonzero Fourier block is $M_\mu\otimes M_\theta$ on the selected
pair of child types, and
$E=\sum_{z=(\alpha,\zeta)\in S_m\times S_m}e(z)z$.
The trace in~\eqref{rec:cutoff-weighted-cycle} is the regular trace
of $S_{2m}$.  To verify it, expand each matching average as a sum over
subsets of the $m$ matched pairs.  For a term to return to the
half-preserving subgroup, its two switch subsets $J,I$ must satisfy
$\alpha(J)=\zeta(J)=I$.  The permissible $J$ are unions of cycles
of $\alpha^{-1}\zeta$, so there are $2^{c(\alpha^{-1}\zeta)}$.
The second half-preserving coefficient exchanges the two maps on $J$,
up to inversion.  This operation is a bijection on the paired terms;
symmetrizing $|e(z)e(z')|\le(|e(z)|^2+|e(z')|^2)/2$ proves the bound.
Schur orthogonality gives
\[
 (m!)^2\sum_z|e(z)|^2=D_\mu t_\mu D_\theta t_\theta.
\]
The cycle factor is therefore averaged under the product of the two
normalized squared-coefficient laws.
Since $(2m)!/4^m\le(m!)^2$, the normalized half-density is
\[
 \phi_\mu(g)=
 \frac{D_\mu|\Tr(M_\mu V_\mu(g))|^2}{t_\mu}.
\]
It is a probability density relative to uniform measure on $S_m$,
by matrix-coefficient orthogonality; we only use it when $t_\mu>0$.
For independent $\alpha,\zeta$ with these respective densities, define
$T_{\mu\theta}=\mathbb E2^{c(\alpha^{-1}\zeta)}$.
The next lemma connects this regular-trace calculation to a fixed parent
representation.  It also identifies the quantities retained when we
later stop a recursion according to its actual traces.

\begin{lemma}[A selected child pair]\label{rec:cutoff-one-node}
Let $B_0,B_1$ be positive operators in the group algebra of $S_m$,
let $B=H(B_0\otimes B_1)H$, and fix $\lambda\vdash2m$ and $p\ge2$.
Let $\mathcal I_\lambda$ consist of the pairs $(\mu,\theta)$ with
$c_{\mu\theta}^\lambda>0$, and write $N_\lambda=|\mathcal I_\lambda|$.
For such a pair use
\[
 M_\mu=B_0(\mu)^{p/2},\quad M_\theta=B_1(\theta)^{p/2},\qquad
 t_\mu=\Tr B_0(\mu)^p,\quad t_\theta=\Tr B_1(\theta)^p
\]
and form $T_{\mu\theta}$ from their coefficient-square densities above.
Pairs with a zero trace are omitted.  Then
\begin{equation}\label{rec:cutoff-one-node-bound}
 D_\lambda\Tr B(\lambda)^p
 \le N_\lambda^2
       \max_{(\mu,\theta)\in\mathcal I_\lambda}
       (D_\mu t_\mu)(D_\theta t_\theta)T_{\mu\theta}.
\end{equation}
If all selected traces vanish, the left side is zero.
\end{lemma}
\begin{proof}
Put $D=B_0\otimes B_1$, and let $E_{\mu\theta}$ be the central
projection onto this child isotypic component.  On $V_\lambda$,
only pairs in $\mathcal I_\lambda$ occur.  The projection includes
all their Littlewood--Richardson multiplicity copies.  Set
$Y_{\mu\theta}=E_{\mu\theta}D^{p/2}$, which is positive because
$E_{\mu\theta}$ commutes with $D$.
Apply \eqref{rec:cutoff-compressed-jensen} on $V_\lambda$, followed
by the Hilbert--Schmidt triangle inequality, to obtain
\[
 D_\lambda\Tr B(\lambda)^p
 \le D_\lambda\left\|\sum_{\mathcal I_\lambda}
                   HY_{\mu\theta}H\right\|_{\HS,V_\lambda}^{2}
 \le N_\lambda^2\max_{\mathcal I_\lambda}
                   D_\lambda\Tr_{V_\lambda}(HY_{\mu\theta}H)^2.
\]
The parent isotypic projection commutes with every group-algebra
operator, including $H$.  Thus the incompatible pairs have been removed
before the next step: the selected positive trace satisfies
\[
 D_\lambda\Tr_{V_\lambda}(HY_{\mu\theta}H)^2
 \le\Tr_{\rm reg,S_{2m}}(HY_{\mu\theta}H)^2.
\]
There is no further restriction-multiplicity factor; those copies were
already included in the parent trace, and the right side includes that
entire trace with regular multiplicity $D_\lambda$.
Cyclicity, \eqref{rec:cutoff-weighted-cycle}, and child orthogonality
bound this regular trace by
\[
 \frac{(2m)!}{4^m(m!)^2}
      (D_\mu t_\mu)(D_\theta t_\theta)T_{\mu\theta}.
\]
The prefactor is at most one.  Notice that the child trace exponent is
still $p$: the matrices in the coefficient densities have exponent $p/2$,
and their Hilbert--Schmidt squares restore $p$.
\end{proof}

The first estimate is representation theoretic.  Put
\[
 G_\mu=\sum_{b=0}^m\ \sum_{\substack{\xi\vdash b\\\zeta\vdash m-b}}
              (c_{\xi\zeta}^\mu)^2.
\]
Here $c_{\xi\zeta}^\mu$ is the Littlewood--Richardson multiplicity of
$V_\xi\otimes V_\zeta$ in the restriction of $V_\mu$ to
$S_b\times S_{m-b}$.  Then
\begin{equation}\label{rec:cutoff-LR-cycle}
 \log_2T_{\mu\theta}\le
               \tfrac12(\log_2G_\mu+\log_2G_\theta).
\end{equation}
The function $2^{c(g)}$ counts the subsets fixed by $g$, so it is
the character of the direct sum of all subset permutation modules.
For an irreducible summand $V_j$ in that sum, its contribution under
the two coefficient-square laws is
$\operatorname{Tr}(\widehat\phi_\mu(j)^*
\widehat\phi_\theta(j))$.  These Fourier matrices are positive
semidefinite, because a square of a coefficient of a positive matrix
is positive definite, and are contractions because the densities
integrate to one.

To bound their traces, let $Q_j$ be the $j$-isotypic projection in
$V_\mu\otimes\overline{V_\mu}$.  Fourier orthogonality gives
\[
 \operatorname{Tr}\widehat\phi_\mu(j)
 =\frac{D_\mu}{t_\mu D_j}
       \operatorname{Tr}[(M_\mu\otimes\overline{M_\mu})Q_j].
\]
Cauchy--Schwarz bounds the last trace by the geometric mean of the
traces with $M_\mu^2$ on either tensor factor.  Each partial trace of
$Q_j$ is scalar by Schur's lemma, so this is at most the multiplicity
of $j$ in $V_\mu\otimes\overline{V_\mu}$ after the displayed
normalization.  Summing over subset modules, Frobenius reciprocity
identifies that sum of multiplicities with
$\sum_b\sum_{\xi,\zeta}(c_{\xi\zeta}^\mu)^2=G_\mu$.
Thus $T_{\mu\theta}$ is at most either $G_\mu$ or $G_\theta$,
and hence at most their geometric mean.  This proves
\eqref{rec:cutoff-LR-cycle}.

Let $u_s(\mu)$ count boxes outside the first $s$ rows and columns,
and put $k_\mu=m-\mu_1$.
With $\eta=10^{-9}$, $\gamma=1/100$ and
$s=\lfloor k_\mu^{1/2-\eta}\rfloor$, one has
\begin{align}
 \log_2G_\mu&\le u_s(\mu)+O(k_\mu^{1-\eta/2}),
                  &&k_\mu\ge m^\gamma,\label{rec:cutoff-core-LR}\\
 \log_2G_\mu&\le k_\mu+O(\log m+\sqrt{k_\mu}\log k_\mu),
                  &&k_\mu<m^\gamma.
                  \label{rec:cutoff-sparse-LR}
\end{align}
For the first, Pieri's rule embeds $V_\mu$ in the induction of its
core together with at most $2s$ one-row or one-column strips.  Restrict
that induced module to two half factors by recording each constituent's
size on each side.  For the core, multiplication of the restriction and
induction dimension bounds gives
\[
 (c_{\xi\zeta}^{\kappa})^2\le\binom{u_s}{|\xi|}\le2^{u_s}.
\]
The strips restrict with multiplicity one, and each subsequent strip
addition also has multiplicity at most one by Pieri.  It remains to count
the intermediate diagrams.  A subdiagram of $\mu$ is specified by its
first row and a partition of at most $k_\mu$ boxes below that row;
its logarithmic count is
$O(\log m+\sqrt{k_\mu}\log(k_\mu+2))$.
Including the size splits, the logarithm of all intermediate choices is
therefore
$O((s+1)(\log m+\sqrt{k_\mu}\log(k_\mu+2)))$.
For $k_\mu\ge m^\gamma$ and
$s=\lfloor k_\mu^{1/2-\eta}\rfloor$, this is
$O(k_\mu^{1-\eta/2})$, after absorbing logarithms.  Squaring and
summing multiplicities changes only the implicit constant and proves
\eqref{rec:cutoff-core-LR}.  For the second, remove only the first row,
using the whole lower diagram as the core in the same argument.
The trivial type costs only $O(\log m)$.

The second cycle estimate depends on the actual trace.  Put
$r_\mu=(\Tr M_\mu)^2/t_\mu$; then
$\phi_\mu\le D_\mu r_\mu$.  The uniform cycle generating product,
split at cycle count $m^{9/10}$ and evaluated at $x=m^{89/100}$,
gives
\begin{equation}\label{rec:cutoff-density-cycle}
 \log_2T_{\mu\theta}\le
 O(1)+m^{9/10}+
 \frac{F_\mu+\log_2r_\mu}{0.88\log_2m}.
\end{equation}
The same bound holds with $\theta$; their average can also be used.
Here is the tail optimization.  For a uniform permutation $g$,
\[
 \mathbb E x^{c(g)}=\prod_{j=0}^{m-1}\frac{x+j}{j+1},\qquad
 \log_2\mathbb E x^{c(g)}=O(x\log m)
 \quad(1\le x\le m).
\]
Take $x=m^{0.89}$, $h=\lceil m^{0.9}\rceil$, and
$A_m=0.88\log_2m$.  Since $m^{0.89}\log m=o(h\log m)$, Markov's
inequality gives, for all sufficiently large $m$,
\[
 \Pr_{\rm unif}\{c(g)\ge j\}\le2^{-A_m j}\qquad(j\ge h).
\]
The law of $\alpha^{-1}\zeta$ has density at most
$B=D_\mu r_\mu$: conditioning on $\zeta$ gives a translate of
the first density, and averaging preserves this bound.  Its cycle tail
is therefore at most $\min(1,B2^{-A_mj})$ for $j\ge h$.
Put $j_*=h+\lceil(\log_2B)/A_m\rceil$.  The tail-sum identity and
$A_m>1$ give
\[
 \mathbb E2^{c(\alpha^{-1}\zeta)}
 \le2^{j_*}+\sum_{j>j_*}2^j B2^{-A_mj}
 \le C2^{j_*}.
\]
Taking logarithms proves \eqref{rec:cutoff-density-cycle}.  The argument
with the two halves interchanged proves its symmetric version.

\subsection{A preliminary trace and the cutoff allocation}

There is a fixed $p_0$ such that, with $\delta=10^{-5}$,
\begin{equation}\label{rec:cutoff-preliminary-trace}
 \Tr K_n^{p_0}|_{V_\lambda}\le2^{\delta F_\lambda}.
\end{equation}
We first specify the stopping contributions.  At a nontrivial sparse
node of size $m$, let $B$ be its positive sweep operator, or any cut
operator satisfying $0\preceq B\preceq T_mT_m^*$.  The local sparse
estimate and $F_\mu=O(k_\mu\log_2m)$ imply
\begin{equation}\label{rec:cutoff-sparse-offset}
 \log_2\bigl(D_\mu\Tr B(\mu)^p\bigr)
       +\tfrac12\log_2G_\mu
 \le 2F_\mu-2p b_1 k_\mu\log_2m
       +\tfrac12\log_2G_\mu\le0
\end{equation}
for one sufficiently large fixed $p$, uniformly over nonzero sparse
nodes.  Here \eqref{rec:cutoff-sparse-LR} bounds the last term by
$O(k_\mu\log_2m)$; level one is annihilated.  Coordinate reversal
permits the same estimate for either positive orientation.  For each
fixed finite set of sizes, the finite-size norm gap likewise permits
one exponent making the left side of \eqref{rec:cutoff-sparse-offset}
nonpositive on every nontrivial type.  A trivial child instead has
$D_\mu=t_\mu=1$ and $G_\mu=m+1$; we stop it and charge its
$O(\log m)$ half-cycle cost to its nontrivial parent.

Choose a large fixed lower size cutoff, and apply
Lemma~\ref{rec:cutoff-one-node} at each remaining nonsparse node,
using the same exponent $p_0$ throughout.  Choose $p_0$ to satisfy
\eqref{rec:cutoff-sparse-offset} and the finite-size stopping inequalities
at this cutoff.  On each selected edge use
\eqref{rec:cutoff-LR-cycle}.  A stopped nontrivial child contributes
its actual $\log_2(D_\mu t_\mu)$ together with its half of this cycle
bound; their sum is nonpositive.  On every active frontier,
Littlewood--Richardson compatibility now gives
\[
 \sum_v k_v\le k,\qquad
 \sum_vF_v\le F_\lambda,\qquad
 \sum_vu_s(v)\le u_s(\lambda).
\]
For the last inequality, consider one Littlewood--Richardson tableau
of $\lambda/\mu$ with content $\theta$.  Entries in the first $s$
rows are at most $s$.  For each value $j>s$, column strictness permits
at most one occurrence in each of the first $s$ columns.  Thus at least
$(\theta_j-s)_+$ occurrences lie beyond both the first $s$ rows and
the first $s$ columns.  These skew core boxes are disjoint from the
$u_s(\mu)$ boxes of the old core, and
$\sum_{j>s}(\theta_j-s)_+=u_s(\theta)$.  Therefore
$u_s(\mu)+u_s(\theta)\le u_s(\lambda)$ at one split; iteration
proves the frontier inequality.
At nodes of size $m$, use the common cutoff
$s_* =\lfloor\max(m^\gamma,\rho m)^\beta\rfloor$,
where $\beta=1/2-3\eta$ and $\rho=k/n$ is the root payload density,
kept fixed throughout this unrolling.
Here is the comparison with the varying cutoff in
\eqref{rec:cutoff-core-LR}.  Put
$a_0=\beta/(1/2-\eta)<1$.  A nonsparse node with its own cutoff
smaller than $s_*$ can occur only when $\rho m\ge m^\gamma$,
and then it has $k_v\le C(\rho m)^{a_0}$.
There are at most $n/m$ nodes of size $m$, so the total level of
these exceptional nonsparse nodes is at most
\[
 C(n/m)(\rho m)^{a_0}
 =Ck(\rho m)^{a_0-1}
 \le Ck m^{-\gamma(1-a_0)}.
\]
Their core counts are at most their levels, so this is a summable error
as the bottom size cutoff increases.  The sparse and finite-size stopping
contributions have already been removed by
\eqref{rec:cutoff-sparse-offset}; trivial-child costs belong to the
active-parent overhead.  All remaining nodes charge at most half their
$u_{s_*}$.
For a parent at level $k_v$, compatible child levels $r,s$ satisfy
$r+s\le k_v$.  Consequently
$N_\lambda\le\sum_{r+s\le k_v}p(r)p(s)$, so the partition-number
bound gives $\log_2N_\lambda=O(\sqrt{k_v}+\log(k_v+2))$.
Restriction overheads and the errors in
\eqref{rec:cutoff-core-LR} are
$O(k_v^{1-\eta/3})$ on nonsparse nodes, so at size $m$ their sum
is $O(km^{-\gamma\eta/3})$.  Summing over dyadic sizes makes their
total $o(k)$ as the bottom cutoff tends to infinity.  These errors
are independent of $p_0$.
A box of index $a$ is charged at at most
$O(1)+L+\beta^{-1}\log_2a$ levels.  By
\eqref{rec:core-charge} the total is at most
$(1/(2\beta)+o(1))F_\lambda<(1+\delta)F_\lambda$.
More explicitly, unrolling \eqref{rec:cutoff-one-node-bound} gives
\[
 F_\lambda+\log_2\Tr K_n(\lambda)^{p_0}
 \le \left(\frac1{2\beta}+\varepsilon(M)\right)F_\lambda,
 \qquad \varepsilon(M)\longrightarrow0
\]
as the fixed lower cutoff $M$ tends to infinity.  The stopping terms are
nonpositive, and the displayed error includes the type counts and
trivial-child costs.  Since $1/(2\beta)<1+\delta/2$, choose $M$ first
so that the coefficient is at most $1+\delta$, and then choose the
single finite exponent $p_0$ as above.  Subtracting $F_\lambda$ proves
\eqref{rec:cutoff-preliminary-trace}.  We may take $M$ still larger
before fixing $p_0$ when a smaller error margin is needed below.

We now turn the preliminary trace into the midpoint projection estimate.
Split the ordered coordinates as evenly as possible, and write
\[
 T_n=C D,
\]
where $D$ applies the first group of coordinates and $C$ the second.
The two coordinate groups have $\lfloor d/2\rfloor$ and
$\lceil d/2\rceil$ coordinates, where $d=\log_2n$.  Each segment
is a product of independent sweeps on its parallel subcubes, of sizes
$2^{\lfloor d/2\rfloor}$ and $2^{\lceil d/2\rceil}$, respectively.
Let $H_D,H_C$ be the corresponding direct product permutation subgroups.
The argument applies equally with the two coordinate-group sizes
interchanged; their middle log sizes differ from $d/2$ by at most $1/2$.
A product type for either subgroup is a tuple of partitions, with product
dimension and hence log dimension equal to the sum of the child log
dimensions.  On $V_\lambda$, let $P_{\rm lo}^D$ and $P_{\rm lo}^C$ be the
central isotypic projections for $H_D$ and $H_C$, respectively, onto
product types whose log dimensions sum to at most
\begin{equation}\label{rec:cutoff-X}
 X=0.50003F_\lambda.
\end{equation}
\begin{proposition}[Two midpoint projections]\label{rec:cutoff-midpoint}
There is an absolute $c_1>0$ such that, for all sufficiently large dyadic
$n$ and all surviving nontrivial $\lambda\vdash n$ with
$k\ge n^{1/100}$,
\begin{equation}\label{rec:cutoff-low-bound}
 \max\left\{\|T_nP_{\rm lo}^D|_{V_\lambda}\|,
              \|P_{\rm lo}^CT_n|_{V_\lambda}\|\right\}
 \le2^{-c_1F_\lambda}.
\end{equation}
The two projections refer to different product subgroups; no commutation
between them is asserted.
\end{proposition}

\begin{proof}
We prove the first estimate; reversing the coordinate order gives the
second.  Consider the positive operator $T_nP_{\rm lo}^DT_n^*$ and take
exponent $p=2p_0$.  Expand its outer matching layers until reaching
the middle coordinate system.  At a node $v$, its operator has the
form $B_v=A_vP_vA_v^*$, where $A_v$ is the remaining local sweep
and $P_v$ is the inherited central cutoff for its innermost product
types.  In particular $0\preceq B_v\preceq A_vA_v^*$.
For a positive contribution on its selected irreducible type define
\[
 M_v=B_v^{p_0},\qquad t_v=\operatorname{Tr}B_v^{2p_0},\qquad
 r_v=(\operatorname{Tr}M_v)^2/t_v.
\]
These are actual cut-operator quantities.  The coefficient-square and
cycle bounds above apply to these positive matrices just as they do to
the unrestricted matrices.
At a binary restriction, allocate its cutoff by the elementary
projection inequality
\begin{equation}\label{rec:cutoff-allocation}
 \mathbf1_{a+b\le B}\le
 \sum_{u+v\le B+2, u,v\in\mathbb Z_{\ge 0}}
          \mathbf1_{a\le u}\otimes\mathbf1_{b\le v}.
\end{equation}
Ceiling $a,b$ proves the cover.  These are commuting central child
projections, so the inequality holds before sandwiching.  If
$\mathcal J_v$ is this finite set of budget pairs, sandwiching gives
\[
 B_v\preceq\sum_{(u,w)\in\mathcal J_v}
 H(B_{0,u}\otimes B_{1,w})H,\qquad
 B_{i,u}=A_iP_{i,u}A_i^*.
\]
The child product log dimensions never exceed $m\log_2m$, so the
ceilings may be capped at $\lceil m\log_2m\rceil$; in particular
$|\mathcal J_v|$ is polynomial in $m$.  Monotonicity of positive
Schatten norms and their triangle inequality, followed by
Lemma~\ref{rec:cutoff-one-node}, give the precise cut recursion
\begin{equation}\label{rec:cutoff-cut-node}
 D_\lambda\Tr B_v(\lambda)^p
 \le |\mathcal J_v|^p N_\lambda^2
 \max_{\substack{(u,w)\in\mathcal J_v\\
                  (\mu,\theta)\in\mathcal I_\lambda}}
      (D_\mu t_{0,u,\mu})(D_\theta t_{1,w,\theta})
                  T_{u,w,\mu,\theta}.
\end{equation}
Here $t_{i,u,\mu}=\Tr B_{i,u}(\mu)^p$, and the cycle factor is
formed from these same actual matrices $B_{i,u}(\mu)^{p/2}$.
Thus allocation adds only $O_{p_0}(\log m)$ to the logarithmic cost
when $p=2p_0$.  The selected cutoff and child pair remain fixed when
we subsequently expand a child trace.

Finally $B_{i,u}\preceq A_iA_i^*$.  Eigenvalue monotonicity bounds
their powered traces by those of the unrestricted operators, whenever
we need the preliminary trace bound or a sparse stopping inequality.
This does not require the power function to be operator monotone.

Active nodes have size at least a constant multiple of $\sqrt n$.
At a fixed size $m$, each nontrivial active node has level at least
$m^{1/100}$, while their total level is at most $k$.  There are therefore
at most $k m^{-1/100}$ such nodes.  The $O_{p_0}(\log m)$ allocation
cost and the bounded rounding loss per split sum to
$O_{p_0}(k m^{-1/100}\log m)$ at that size.  Summing over the dyadic
sizes down to the middle gives $o(k)$.  Trivial children are stopped and
their costs have already been assigned to the parent overhead.  The
terminal sum includes only branches reaching the middle: stopped
nontrivial branches have already contributed nonpositive terms by
\eqref{rec:cutoff-sparse-offset}.  The successive ceilings therefore give
$\sum_vF_v\le X+o(k)$ at this terminal cut, and
\begin{equation}\label{rec:cutoff-terminal-budget}
 \sum_v\log_2(D_vt_v)\le(1+\delta)(X+o(k))
\end{equation}
by \eqref{rec:cutoff-preliminary-trace}.

\subsection{The actual-trace invariant}

Before estimating the total cycle cost, we specify a second stopping rule.
Put $d=\log_2n$.  While descending through child log sizes between
$d$ and $0.9d$, test the actual frontier traces against
\begin{equation}\label{rec:cutoff-diagnostic}
                 -\sum_v\log_2t_v>0.2F_\lambda.
\end{equation}
If this inequality holds, the trace already saved at that frontier will
replace the more expensive terminal budget at the middle.  The following
invariant explains why this decision can be made using the actual traces.

A frontier consists of the nodes at the current common coordinate depth,
with previously stopped branches removed and their paid contributions
retained separately.  Expand every currently active node once to form a
complete new generation.  Test that generation before removing any of
its newly sparse or trivial children.  At each expansion choose a
maximizing term in
\eqref{rec:cutoff-cut-node}, using its actual child traces and actual
weighted-cycle factor.  Retain the selected matrices and the resulting
number $T_e$.  Subsequent expansion replaces only a child quantity
$q_v=\log_2(D_vt_v)$, without changing that earlier density.

To account explicitly for stopped children, start an edge cost at
$\Gamma_e=\log_2T_e$.  When its child $w$ is stopped, write
\[
 q_w+\Gamma_e=
 \left(q_w+\tfrac12\log_2G_w\right)
 +\left(\Gamma_e-\tfrac12\log_2G_w\right).
\]
The first term is nonpositive for a nontrivial stopped child by
\eqref{rec:cutoff-sparse-offset}; for a trivial child it is the
$O(\log m)$ parent overhead already recorded.  Keep the second term
as the new residual cost $\Gamma_e$.  This scalar subtraction does not
change the actual matrix or coefficient density used to define $T_e$.
After this stopping decision, freeze $\Gamma_e$ as well; expanding an
unstopped child later changes only that child's frontier contribution.
The two usable bounds are consequently
\[
 \Gamma_e\le\log_2T_e,\qquad
 \Gamma_e\le\tfrac12\sum_{w\text{ unstopped child of }e}\log_2G_w.
\]
The first permits the density estimate, and the second permits the core
charge without charging stopped children again.
At every intermediate frontier,
\begin{equation}\label{rec:cutoff-actual-invariant}
 \log_2(D_\lambda t_\lambda)
 \le \sum_{e\ \rm retained}\Gamma_e+
       \mathrm{overheads}+\sum_{v\ \rm frontier}\log_2(D_vt_v)
       +\mathrm{paid\ nontrivial\ stopping\ terms}.
\end{equation}
The last terms are nonpositive.  This follows inductively from
\eqref{rec:cutoff-cut-node} and the displayed scalar decomposition.
In particular no independence between later maximizing choices is needed,
and a newly created frontier is tested before any of its offsets are
subtracted.

It remains to show that the resulting cost leaves fixed slack below the
root dimension.  We first estimate what the diagram-only cycle bound would
cost throughout the truncated recursion.  Put $l=L/d$ and
$z=\log_2a/d$ for a charged box.
Over the relevant sizes, from $d/2$ to $d$, its core charge is at most
\begin{equation}\label{rec:cutoff-box-cost}
 O(1)+\tfrac d2[\min(1,l+\beta^{-1}z)-1/2]_+.
\end{equation}
We charge only boxes appearing in at least one of these cores.  In
particular $l+\beta^{-1}z\ge1/2-O(1/d)$, so their weights
$w_x=d(l+z)$ are at least $\beta d/2-O(1)$.  The $O(1)$ rounding
terms therefore total $O(F_\lambda/d)$ by \eqref{rec:core-charge}.
For the principal term put $q=\beta^{-1}$ and $y=l+qz$.  Its ratio
to $w_x$ satisfies
\[
 \frac{[\min(1,y)-1/2]_+}{2(l+z)}
 \le\frac q4<0.50001,
\]
because $l+z\ge y/q$.  If $z<0.49$, set $b_z=(q-1)z<1/2$;
the same ratio equals
$[\min(1,y)-1/2]_+/[2(y-b_z)]$.  It increases on $1/2<y<1$
and decreases on $y>1$, hence is at most
\[
 \frac1{4(1-b_z)}\le
 \frac1{4(1-0.49(q-1))}<0.491.
\]
These bounds hold on the whole range $l,z\ge0$, and therefore on the
admissible diagram range.  The additional geometric relation
$a(x)(a(x)-1)\le k$ will be used below to turn a deep box into a
lower bound for $k$.

Choose the fixed cutoff in the preliminary proof so that its
core-charge and summable-error remainders are sufficiently small, and
then fix $p_0$.  Increasing the root size afterward makes the allocation
and rounding errors small as well.  We take the sum of the error terms
in the following budget comparisons to be at most $10^{-5}F_\lambda$;
the displayed numerical comparisons all leave more slack than this.
If the total core charge is at most
$0.4995F_\lambda$, combining it with
\eqref{rec:cutoff-terminal-budget} suffices, since
\begin{equation}\label{rec:cutoff-ordinary-slack}
 0.4995+(1.00001)(0.50003)=0.9995350003<1.
\end{equation}

Otherwise split the dimension budget, rather than the incurred cycle
charge, according to the value of $z$.  Give a charged box the weight
$w_x=d(l+z)=L+\log_2a(x)$, and let $W_{\rm hi},W_{\rm lo}$ be the
sums of these weights over $z\ge0.49$ and $z<0.49$, respectively.
The core-charge estimate gives
$W_{\rm hi}+W_{\rm lo}\le(1+o(1))F_\lambda$; the omitted shallow
boxes and the other error terms cost $o(F_\lambda)$ after the fixed
bottom cutoff has been chosen large enough.  The two ratios following
\eqref{rec:cutoff-box-cost} therefore imply
\[
 0.4995F_\lambda
 <0.50001W_{\rm hi}+0.491W_{\rm lo}+o(F_\lambda)
 \le0.491F_\lambda+0.00901W_{\rm hi}+o(F_\lambda).
\]
Since $0.0085/0.00901>0.943$, it follows, with room for those errors,
that $W_{\rm hi}\ge0.93F_\lambda$.  In particular some box has
$a\ge n^{0.49}$, so the diagram contains at least $a(a-1)$ boxes
below its first row and $k\ge n^{0.97}$ for large $n$.
In the lower portion, from log size $d/2$ to $0.9d$, each box costs
at most $0.2d$.  Its cost-to-budget ratio on the high set is thus at
most $0.2/0.49$, while on the low set the full bound $0.50001$ still
applies.  The total lower cost is at most
\[
 \frac{0.2}{0.49}W_{\rm hi}+0.50001W_{\rm lo}+o(F_\lambda)
 \le\left\{0.93\frac{0.2}{0.49}+0.07(0.50001)+o(1)\right\}F_\lambda
 <0.418F_\lambda.
\]
In the upper portion use \eqref{rec:cutoff-density-cycle}, testing
\eqref{rec:cutoff-diagnostic} at each new frontier as specified above.

As long as \eqref{rec:cutoff-diagnostic} has not fired,
\eqref{rec:cutoff-preliminary-trace} and monotonicity for cut
operators give
\[
 \sum_v\log_2r_v
 =2\sum_v\log_2\Tr M_v-\sum_v\log_2t_v
 \le(0.2+2\delta)F_\lambda.
\]
At size $m\ge n^{0.9}/2$, the additive $m^{0.9}$ terms in
\eqref{rec:cutoff-density-cycle}, summed over at most $n/m$ nodes,
are $O(nm^{-0.1})=O(n^{0.91})$.  Summing over these dyadic sizes
keeps this bound, which is $o(F_\lambda)$ because
$F_\lambda\ge b n^{0.97}$.  The type and allocation overheads are
also $o(F_\lambda)$ by the previous bounds.
The remaining inverse-log-size sum gives an upper-portion cost of at most
\begin{equation}\label{rec:cutoff-upper-cost}
 \frac{1.2+2\delta}{2(0.88)}\ln(1/0.9)F_\lambda+o(F_\lambda)
 <0.076F_\lambda.
\end{equation}
The factor $1/2$ is the averaging of the two child bounds; the
logarithm is the integral of inverse log size across this portion.
If the diagnostic never fires, the total charge is below
$(0.418+0.076)F_\lambda=0.494F_\lambda$.
If it fires, test it on the newly created frontier before removing any
newly sparse children.  Previously stopped sparse nodes have already
paid their edge offsets and retain nonpositive contributions.  At the
triggering frontier the actual sum
$\sum_v\log_2(D_vt_v)$ is below $0.8F_\lambda$.
Use~\eqref{rec:cutoff-actual-invariant} immediately with that sum.
The density-cycle estimate cannot be charged on this triggering level,
because its diagnostic hypothesis has failed there.  Use the crude
Littlewood--Richardson bound instead.  Its core part charges boxes with
$\log_2a$ of order $d$, so~\eqref{rec:core-charge} bounds the
single-level charge by $O(F_\lambda/d)=o(F_\lambda)$.
For newly sparse children, the additional bound in
\eqref{rec:cutoff-sparse-LR} is also $o(F_\lambda)$, since there
are at most $n/m$ children of size $m\ge2^{.9d-O(1)}$ and
$F_\lambda\ge b n^{.97}$.  The total is therefore below
$0.876F_\lambda+o(F_\lambda)$.  Both cases leave fixed slack below one.
In every case we have
$\log_2(D_\lambda t_\lambda)\le(1-\varepsilon_0)F_\lambda$
for some fixed $\varepsilon_0>0$, at all sufficiently large active
roots.  Thus $t_\lambda\le2^{-\varepsilon_0F_\lambda}$.
Since $T_nP_{\rm lo}^DT_n^*$ has norm
$\|T_nP_{\rm lo}^D\|^2$ and $t_\lambda$ is its $2p_0$ moment,
taking a $4p_0$-th root proves~\eqref{rec:cutoff-low-bound}.
Its constants and size threshold are independent of the exponent used
in the subsequent norm induction.
\end{proof}

\subsection{Closing the norm induction without a circular base}

Write $P_D=P_{\rm lo}^D$ and $P_C=P_{\rm lo}^C$.  With $T_n=CD$,
\[
 T_n=T_nP_D+P_CT_n(I-P_D)+(I-P_C)CD(I-P_D).
\]
Proposition~\ref{rec:cutoff-midpoint} bounds the first two terms by
$2^{-c_1F_\lambda}$ each.  The central projections commute with their
own segments, $P_DD=DP_D$ and $P_CC=CP_C$.  On a high product type,
the child dimension-power estimates multiply, and its product log
dimension exceeds $X$.  Thus, if all child estimates hold with exponent
$g$, the last term has norm at most
\[
 \|(I-P_C)C\|\,\|D(I-P_D)\|
 \le2^{-2gX}=2^{-1.00006gF_\lambda}.
\]
This uses no commutation between $P_C$ and $P_D$.
We need a size threshold independent of the final small exponent.
A crude polynomial gap supplies it.  For a unit vector $v$, put
$\varepsilon=1-\|T_nv\|^2$.  If $v_i$ is the vector after $i$
matching projections, orthogonality gives
$\sum_i\|v_i-v_{i-1}\|^2=\varepsilon$.
Hence $v$ is within $\sqrt{d\varepsilon}$ of every matching-fixed
space.  Each edge transposition in that matching moves $v$ by at most
$2\sqrt{d\varepsilon}$.  Any transposition is a product of at most
$2d-1$ cube-edge transpositions along a cube path, and any permutation
is a product of at most $n-1$ arbitrary transpositions.  Telescoping
therefore bounds $\|\pi v-v\|$ by $Cnd\sqrt{d\varepsilon}$.
The average of $\pi v$ over $S_n$ is zero on a nontrivial irreducible,
so $1\le Cnd\sqrt{d\varepsilon}$.
Together with $F_\lambda\le n\log_2n$, this supplies an exponent
$a n^{-C_0}$ in the bound $\|T_n\|\le2^{-gF_\lambda}$,
for absolute $a,C_0>0$; dimension-one nontrivial types are annihilated.

Fix temporarily a desired exponent $0<c\le c_1/2$.  At a parent of
size $n$, induction and this independent gap permit a common child
exponent
\[
 g=\max\{c,\min(c_1/2,a n^{-C_0})\}.
\]
Thus $g\le c_1/2$ and $g\ge a'n^{-C_0}$ for an absolute $a'>0$.
Relative to $2^{-gF_\lambda}$, the saving in the high--high term is
$1-2^{-.00006gF_\lambda}$.  It is bounded below by an inverse
polynomial in $n$, whereas the two low terms have relative size at most
$2\cdot2^{-(c_1/2)F_\lambda}$.  On active roots,
$F_\lambda\ge b n^{.01}$, so one absolute large-size threshold
absorbs the low terms for every such $c$.  Below that threshold choose
$c>0$ small enough to fit every finite nontrivial norm gap, and also
small enough for~\eqref{rec:core-dimension-sparse} on sparse roots,
where $F_\lambda=O(k\log n)$.  This closes the induction without
a circular choice of its base and proves
Theorem~\ref{rec:core-power-main}.

For completeness, let $\mu_q$ be the permutation law after $q$
sweeps and define its squared relative $L^2$ deviation by
$\chi^2=n!\sum_{g\in S_n}|\mu_q(g)-1/n!|^2$.
The regular Fourier consequence is
\begin{equation}\label{rec:cutoff-fourier-completion}
 \chi^2\le\sum_{\lambda\ne(n),\ T_n(\lambda)\ne0}
             D_\lambda^2\,2^{-2qcF_\lambda}=o(1)
\end{equation}
for a sufficiently large absolute $q$.  There are at most
$2^{O(\sqrt k\log(k+2))}$ types at level $k$, $F_\lambda\ge bk$,
and $F_\lambda\to\infty$ at each fixed positive $k$.
These facts justify the last limit and give the same bound from
every starting deck by translation.
